\documentclass[11pt]{article}
\usepackage{preamble}
\setlength{\parskip}{4pt}

\begin{document}

\daybanner{AP Statistics \textperiodcentered\ Unit 4 \textperiodcentered\ Topic 4.6 \textperiodcentered\ B: 2027-01-22 \textperiodcentered\ E: 2027-03-10 \textperiodcentered\ TEACHER KEY}{Two Samples, One Difference}

\sectionbanner{CED TETHER}

{\small
\begin{itemize}[leftmargin=1.5em,itemsep=2pt,topsep=2pt]
  \item \textbf{ENDURING UNDERSTANDING: UNC-3} --- Probabilistic reasoning allows us to anticipate patterns in data.
  \item \textbf{Skills: 3.B} --- Determine parameters for probability distributions. | 3.C --- Describe probability distributions. | 4.B --- Interpret statistical calculations and findings to assign meaning or assess a claim.
  \item UNC-3.T Determine parameters of a sampling distribution for a difference in sample means. [Skill 3.B]
  \item UNC-3.T.1 For a numerical variable, when randomly sampling from two independent populations with means $\mu$1 and $\mu$2 and standard deviations $\sigma$1 and $\sigma$2, the sampling distribution of the difference in sample means has mean $\mu$ = $\mu$1 - $\mu$2 and standard deviation $\sigma$ = sqrt($\sigma$1\textsuperscript{2}/n1 + $\sigma$2\textsuperscript{2}/n2). UNC-3.T.2 If sampling without replacement and sample sizes are less than 10\% of population sizes, the difference is negligible.
  \item UNC-3.U Determine whether a sampling distribution of a difference in sample means can be described as approximately normal. [Skill 3.C]
  \item UNC-3.U.1 The sampling distribution of the difference in sample means can be modeled with a normal distribution if both population distributions can be modeled with normal distributions. UNC-3.U.2 If population distributions cannot be modeled with normal distributions, the sampling distribution can be approximated by a normal distribution if both sample sizes are at least 30.
  \item UNC-3.V Interpret probabilities and parameters for a sampling distribution for a difference in sample means. [Skill 4.B]
  \item UNC-3.V.1 Probabilities and parameters for a sampling distribution for a difference of sample means should be interpreted using appropriate units and within the context of specific populations.
\end{itemize}
}
\textbf{Essential question:} How can two sample means use different pathways to justify one normal model for their difference?\par
\textbf{DOK-3 skill:} 4.B \textperiodcentered\ \textbf{FRQ pattern:} {\small\texttt{two-\allowbreak{}shape-\allowbreak{}routes-\allowbreak{}for-\allowbreak{}a-\allowbreak{}mean-\allowbreak{}difference}}\par
\textbf{DOK rationale:} Part (c) coordinates a separate shape pathway for each sample mean, judges two incomplete model claims, and uses a small-sample counterfactual to trace when the tail probability loses support.\par

\frameworkphaseheader{Do Now (first take) $\rightarrow$ Explore (video + follow-along) $\rightarrow$ Exit (finish + turn in)}{1 $\rightarrow$ 3}{45}{%
\begin{itemize}[leftmargin=1.5em,itemsep=2pt,topsep=2pt]
  \item Board slide up at the bell. Ask which details govern the sampling model.
  \item Begin the lesson video follow-along at minute 5.
  \item Return at minute 33. Require separate shape arguments and the counterfactual check.
  \item Collect at the door. Score part (c) only, E/P/I.
\end{itemize}
}{%
\begin{itemize}[leftmargin=1.5em,itemsep=2pt,topsep=2pt]
  \item Name one argument to retain and one that needs more evidence.
  \item Complete the follow-along, finish (a)--(c), revise, and turn in.
\end{itemize}
}{%
\begin{itemize}[leftmargin=1.5em,itemsep=2pt,topsep=2pt]
  \item Which distribution are we justifying for North: individual waits or sample means?
  \item Why can the two groups use different routes to approximate normality?
  \item What changes if the skewed South sample has size 8?
  \item What repeated-sampling event does 0.0273 describe?
\end{itemize}
}{Solo teacher. Justify North and South separately, then combine them using independence.}

\begin{callout}{\IconWarn\ WATCH FOR}{calloutred}
\begin{itemize}[leftmargin=1.5em,itemsep=2pt,topsep=2pt]
  \item Requiring both samples to have size at least 30 even when one population is normal.
  \item Assuming one normal component automatically fixes an unjustified second component.
  \item Subtracting variances because the statistic subtracts means.
\end{itemize}
\end{callout}

\sectionbanner{SCORING PART (c) --- E / P / I}

\textbf{Expected elements}
\begin{itemize}[leftmargin=1.5em,itemsep=2pt,topsep=2pt]
  \item \textbf{judgment-1} (required) --- Supports the original probability using North population normality and South n=48 CLT route.
  \item \textbf{judgment-2} (required) --- Uses independence and explains the incomplete reasoning in both accounts.
  \item \textbf{judgment-3} (required) --- Uses the 48-plus-8 counterexample: total 56 does not justify the small skewed South component.
\end{itemize}

{\small\renewcommand{\arraystretch}{1.25}
\noindent\begin{tabular}{|p{0.5in}|p{5.9in}|}
\hline
\textbf{E} & Meets all three checks with correct contextual reasoning; equivalent wording is acceptable. \\ \hline
\textbf{P} & Shows at least one substantive correct check, but has an omission or error that prevents E. \\ \hline
\textbf{I} & Shows none of the three checks with substantive correct reasoning; a choice alone is insufficient. \\ \hline
\end{tabular}}

\clearpage
\focusbanner{TODAY'S PROBLEM --- annotated}

Suppose parcel pickup waits have the shown population parameters, measured in minutes. There are 600 North visits and 1,000 South visits. North waits are approximately normal; South waits are strongly right-skewed. Independent SRSs are taken without replacement, and the service wants $P(\bar X_N-\bar X_S>2.5)$.\par\smallskip \textbf{Rina:} ``South is skewed, so I would not use a normal model for the difference.''\par \textbf{Oscar:} ``North is normal and the combined sample size is 60, so the normal model is valid; the probability is $0.0273$.''\par
\begin{center}
\textbf{Locker samples}\par\smallskip
\begin{tabular}{|l|c|c|c|}
\hline
\textbf{Locker} & \textbf{$n$} & \textbf{$\mu$} & \textbf{$\sigma$} \\ \hline
\textbf{North} & 12 & 6.0 & 1.8 \\ \hline
\textbf{South} & 48 & 4.8 & 3.0 \\ \hline
\end{tabular}
\end{center}
\begin{firsttakebox}{FIRST TAKE (5 min)}
Would you check each place's sample size separately or add them together? Give one reason from the study.\par
\answer{Not graded. Each account contains a claim that still needs component-level justification.}
\end{firsttakebox}

\textit{Rules callout ``JUSTIFY EACH SAMPLE MEAN BEFORE THEIR DIFFERENCE (keep on the page)'' is printed on the student sheet.}\par\medskip

\dokbadge{1}~\textbf{(a)} State the mean of $\bar X_N-\bar X_S$ and translate the event above into context.\par
\answer{The sampling distribution is centered at $\mu_N-\mu_S=6.0-4.8=1.2$ minutes. The event says the sample mean North wait exceeds the sample mean South wait by more than 2.5 minutes.}

\dokbadge{2}~\textbf{(b)} Verify both 10 percent conditions, compute the sampling-distribution SD, and calculate the probability under a normal model.\par
\answer{North: $12\leq0.10(600)=60$; South: $48\leq0.10(1{,}000)=100$. The SD is $\sqrt{1.8^2/12+3.0^2/48}=\sqrt{0.4575}=0.676387$. Thus $z=(2.5-1.2)/0.676387=1.92198$ and the upper-tail probability is $0.027304\approx0.0273$.}

\dokbadge{3}~\textbf{(c)} In 4--5 sentences, decide whether 0.0273 is supported and explain what each student overlooks. Justify each sample mean using its own shape and size. Then consider $n_N=48$ and $n_S=8$: explain why their total alone cannot justify the same normal-model claim.\par
\answer{The original probability is supported, but both explanations are incomplete. North's normal population supports its sample mean at $n_N=12$, and South's $n_S=48$ supports its sample mean by the CLT. Independence then supports an approximately normal difference, so Rina rejects a supported approximation. Oscar's combined-size rule fails: with $n_N=48$ and $n_S=8$, the total is 56, but the small sample from skewed South has no normality justification. The total alone would make readers trust an unsupported tail approximation.}

\begin{summaryexitbox}{EXIT}
One thing the video changed about my first take: \hrulefill
\end{summaryexitbox}

\end{document}
